A elicitação de requisitos de usabilidade permanece um desafio recorrente na Engenharia de
Software, especialmente pela dificuldade de tornar explícitos aspectos de usabilidade ainda nas
fases iniciais do desenvolvimento. O método USARP (Usability Requirements with Personas
and User Stories) foi proposto com o objetivo de apoiar a identificação, discussão e especificação
sistemática de requisitos de usabilidade por meio do uso integrado de personas, user stories e
mecanismos de usabilidade representados por cartas. Apesar de estudos empíricos indicarem
benefícios associados ao método, ainda carece de uma síntese que integre as evidências disponíveis
e avalie sua aplicação em diferentes contextos. Este trabalho apresenta uma síntese empírica do
método USARP, fundamentada na integração de evidências quantitativas e qualitativas provenientes
de quatro estudos primários conduzidos em contextos acadêmicos e industriais. A metodologia
combina análise estatística descritiva, redução de dimensionalidade por meio da Análise de
Componentes Principais (PCA) e análise temática das respostas abertas, consolidadas no nível do
respondente, permitindo investigar dimensões latentes relacionadas ao valor percebido do método
e à utilidade percebida de seus mecanismos. Os resultados indicam percepção amplamente positiva
e homogênea do método entre os contextos acadêmico e industrial, com forte efeito de teto nos
itens avaliados. A Análise de Componentes Principais, de caráter exploratório, evidenciou uma
dimensão de percepção positiva geral e contrastes secundários entre os blocos de cartas do método.
As comparações entre contextos não revelaram diferenças estatisticamente robustas após a correção
para múltiplas comparações, com exceção de um item relativo ao contexto do mecanismo de
usabilidade. A análise qualitativa das respostas abertas identificou as cartas e os guias de
especificação como o principal elemento de apoio à elicitação, além de tensões relacionadas à curva
de aprendizado e ao volume de cartas e da subutilização das cartas de prototipação. Como
contribuição, este trabalho fornece uma síntese empírica crítica do USARP, destacando suas
potencialidades, limitações e implicações para sua adoção em contextos reais de desenvolvimento
de software.

\palavraschave{USARP; requisitos de usabilidade; síntese empírica; engenharia de software; personas e user stories.}